% 全书科研图字体、语义颜色与线型的统一覆盖层。
% 参考 lcpmgh/colors @palettes.csv 中 id 23/29/36/53 与图3.14。
\renewcommand{\figurefont}{\sffamily\CJKfamily{bookkai}}
\colorlet{FigureInput}{FigureBlue}
\colorlet{FigureModel}{FigureYellow}
\colorlet{FigureOutput}{FigureRed}
\colorlet{FigureLoss}{FigureRed}
\definecolor{FigureStroke}{HTML}{4C4D4F}
\definecolor{FigureMuted}{HTML}{747A80}
\definecolor{FigureGrid}{HTML}{C8CDCF}
\colorlet{FigureCream}{FigureYellowFill}
\colorlet{FigureInputFill}{FigureBlueFill}
\definecolor{FigureModelFill}{HTML}{FFFFFF}
\colorlet{FigureOutputFill}{FigureRedFill}
\colorlet{FigureLossFill}{FigureRedFill}
\definecolor{FigurePanel}{HTML}{F4F5F5}
% xcolor 的兼容别名是颜色副本，覆盖基础色后必须一并重建。
\colorlet{NNInk}{FigureInk}
\colorlet{NNInput}{FigureBlue}
\colorlet{NNHidden}{FigureRed}
\colorlet{NNOutput}{FigureYellow}
\colorlet{NNGradient}{FigureLoss}
\colorlet{NNLine}{FigureMuted}
\colorlet{NNPanel}{FigurePanel}
\colorlet{NNWire}{FigureStroke}
\colorlet{FoundationInput}{FigureInput}
\colorlet{FoundationModel}{FigureModel}
\colorlet{FoundationOutput}{FigureOutput}
\colorlet{FoundationLoss}{FigureLoss}
\pgfplotsset{every axis/.append style={
  font=\figurefont\footnotesize,
  tick label style={font=\figurefont\footnotesize,text=FigureInk},
  label style={font=\figurefont\footnotesize,text=FigureInk},
  title style={font=\figurefont\bfseries\footnotesize,text=FigureInk},
  axis line style={draw=FigureStroke,line width=.6pt},
  tick style={draw=FigureStroke,line width=.6pt},
  grid style={draw=FigureGrid,line width=.5pt},
  legend style={font=\figurefont\footnotesize,text=FigureInk,draw=none,fill=none}
}}
\tikzset{figure heading/.append style={font=\figurefont\mdseries\footnotesize}}
\pgfplotsset{every axis/.append style={
  title style={font=\figurefont\mdseries\footnotesize}}}

\tikzset{book node/.style={figure node,minimum height=10mm,inner sep=6pt,font=\figurefont\small},book arrow/.style={figure flow}}
